\section{Metodologia}

\subsection{Base de dados e desfecho}

Foram utilizadas as 400 imagens e as 11.534 anotações celulares da CRIC Cervix
\cite{rezende2021cric}. A classe normal reúne 6.779 células NILM. A classe
anormal reúne 4.755 células: 606 ASC-US, 1.360 LSIL, 925 ASC-H, 1.703 HSIL e
161 SCC, segundo o Sistema Bethesda \cite{nayar2015bethesda}. A escolha binária
representa um primeiro estágio de priorização e não substitui a classificação
Bethesda.

Para cada anotação, foi gerado um recorte RGB de $224\times224$ pixels centrado
na coordenada do núcleo. Recortes que ultrapassavam as bordas foram completados
com fundo branco. Não foi realizada segmentação do núcleo ou do citoplasma,
permitindo que a rede utilizasse o contexto celular local. A
Figura~\ref{fig:amostras} ilustra recortes reais utilizados no treinamento e na
avaliação.

\begin{figure}[ht]
\centering
\includegraphics[width=.82\textwidth]{amostras_cric_cv.png}
\caption{Exemplos de recortes reais da CRIC Cervix usados no protocolo binário.
A linha superior mostra células NILM; a inferior inclui categorias Bethesda
anormais.}
\label{fig:amostras}
\end{figure}

\subsection{Validação cruzada agrupada}

A unidade de agrupamento foi a imagem original. A avaliação empregou
\texttt{StratifiedGroupKFold} com cinco folds, usando a imagem-fonte como grupo
e preservando, tanto quanto possível, a proporção das classes em cada divisão.
Em cada fold, todas as células de uma mesma imagem permanecem juntas: ou no
conjunto de treino do fold, ou no conjunto de avaliação do fold. Essa decisão
reduz vazamento por correlação visual entre recortes da mesma imagem.
As divisões foram geradas com semente fixa (42) e reutilizadas para o modelo
principal e para o comparador, garantindo que diferenças observadas não fossem
causadas por sorteios distintos dos folds. A comparação é, portanto, pareada no
nível do fold.

A ConvNeXt-Tiny foi o modelo principal do estudo. Para comparação complementar,
uma ResNet-50 foi treinada e avaliada nos mesmos folds. Dentro do conjunto de
treino de cada fold, um \texttt{GroupShuffleSplit} separou uma validação interna
agrupada por imagem, usada somente para selecionar o melhor checkpoint pela AUC
ROC. O conjunto de avaliação do fold não participou dessa seleção. As predições
foram avaliadas no limiar 0,50 e sem \emph{test-time augmentation}. A
Figura~\ref{fig:pipeline} resume o fluxo experimental.

\begin{figure}[ht]
\centering
\resizebox{\textwidth}{!}{%
\begin{tikzpicture}[
  node distance=0.35cm,
  etapa/.style={
    draw,
    rounded corners=2pt,
    align=center,
    font=\scriptsize,
    minimum height=0.70cm,
    text width=1.70cm,
    inner sep=2pt
  },
  seta/.style={-{Latex[length=1.4mm]},thick}
]
\node[etapa] (img) {Imagem\\CRIC};
\node[etapa,right=of img] (crop) {Recortes\\$224\times224$};
\node[etapa,right=of crop] (fold) {5 folds\\por imagem};
\node[etapa,right=of fold] (train) {Ajuste fino\\por fold};
\node[etapa,right=of train] (models) {ConvNeXt\\ResNet};
\node[etapa,right=of models] (eval) {Avaliação\\do fold};
\node[etapa,right=of eval] (cmp) {Comparação\\pareada};
\draw[seta] (img) -- (crop);
\draw[seta] (crop) -- (fold);
\draw[seta] (fold) -- (train);
\draw[seta] (train) -- (models);
\draw[seta] (models) -- (eval);
\draw[seta] (eval) -- (cmp);
\end{tikzpicture}
}
\caption{Fluxo do protocolo: a ConvNeXt-Tiny é avaliada em validação cruzada
agrupada por imagem, e a ResNet-50 é usada como comparador complementar nos
mesmos folds.}
\label{fig:pipeline}
\end{figure}

\subsection{Modelos e treinamento}

A arquitetura proposta foi ConvNeXt-Tiny \cite{liu2022convnext}, pré-treinada
na ImageNet \cite{deng2009imagenet} e com a cabeça substituída por uma camada
linear de duas saídas. A ResNet-50 \cite{he2016resnet} recebeu a mesma adaptação
e foi usada apenas como referência arquitetural. O ajuste fino utilizou AdamW
\cite{loshchilov2019adamw}, taxa inicial $3\times10^{-4}$, decaimento
$10^{-4}$, \emph{label smoothing} 0,03, lote 48 e agendador cossenoidal. No
treino, foram aplicados recorte redimensionado aleatório, espelhamento
horizontal, rotação de até 10 graus e pequenas variações de cor; na avaliação,
apenas redimensionamento e normalização pelas estatísticas da ImageNet.

Os folds da ConvNeXt-Tiny haviam sido ajustados inicialmente com teto de 12
épocas. A análise dos históricos mostrou que o melhor checkpoint por AUC de
validação ocorreu nas épocas 1, 3, 3, 4 e 6, conforme o fold. Portanto, a
ResNet-50 foi ajustada com teto de 6 épocas e parada antecipada de paciência 2,
mantendo o mesmo critério de seleção. Essa decisão reduz custo computacional e
permanece dentro da janela efetiva em que a ConvNeXt-Tiny atingiu seus melhores
checkpoints; assim, a comparação prioriza o desempenho no mesmo particionamento,
e não o custo de treinamento das arquiteturas.

Todos os artefatos foram preservados por fold: definição das imagens em cada
divisão, histórico de treinamento, checkpoint selecionado e métricas do conjunto
de avaliação. O treinamento foi executado em GPU NVIDIA RTX 3060; a
ConvNeXt-Tiny e a ResNet-50 possuem cerca de 28 e 25 milhões de parâmetros,
respectivamente. Esse registro favorece reprodutibilidade e segue recomendações de
relato para estudos de modelos preditivos aplicados à saúde
\cite{collins2024tripod,tejani2024claim}.

\subsection{Métricas}

Foram calculadas AUC ROC, AUC precisão--revocação (AUC PR), sensibilidade,
especificidade, acurácia balanceada, \emph{F1} da classe anormal, coeficiente de
correlação de Matthews (MCC), kappa de Cohen e \emph{Brier score}. Todas as
métricas foram calculadas separadamente no conjunto de avaliação de cada fold e
resumidas como média e desvio-padrão.
